\documentclass[10pt,a4paper]{amsart}
\usepackage{amsmath,amssymb,amsthm,graphicx}
% Retained minimal preamble for the Stacks Project chapter varieties.tex.
% Verbatim excerpt of the chapter preamble: Xy-pic setup (native preamble lines 14-18), the native theorem
% declarations with the mathematical macros (native preamble lines 154-217), and the project own no-op
% definitions of the reference/slogan/history display environments (native wrapper.tex lines 1-3) that the
% counted statement uses for its slogan. Copied mechanically; nothing added, removed or redefined.
% For commutative diagrams we use Xy-pic
\usepackage[all]{xy}

% We use 2cell for 2-commutative diagrams.
\xyoption{2cell}
% Theorem environments.
%
\theoremstyle{plain}
\newtheorem{theorem}[subsection]{Theorem}
\newtheorem{proposition}[subsection]{Proposition}
\newtheorem{lemma}[subsection]{Lemma}

\theoremstyle{definition}
\newtheorem{definition}[subsection]{Definition}
\newtheorem{example}[subsection]{Example}
\newtheorem{exercise}[subsection]{Exercise}
\newtheorem{situation}[subsection]{Situation}

\theoremstyle{remark}
\newtheorem{remark}[subsection]{Remark}
\newtheorem{remarks}[subsection]{Remarks}

\numberwithin{equation}{subsection}

% Macros
%
\def\lim{\mathop{\mathrm{lim}}\nolimits}
\def\colim{\mathop{\mathrm{colim}}\nolimits}
\def\Spec{\mathop{\mathrm{Spec}}}
\def\Hom{\mathop{\mathrm{Hom}}\nolimits}
\def\Ext{\mathop{\mathrm{Ext}}\nolimits}
\def\SheafHom{\mathop{\mathcal{H}\!\mathit{om}}\nolimits}
\def\SheafExt{\mathop{\mathcal{E}\!\mathit{xt}}\nolimits}
\def\Sch{\mathit{Sch}}
\def\Mor{\mathop{\mathrm{Mor}}\nolimits}
\def\Ob{\mathop{\mathrm{Ob}}\nolimits}
\def\Sh{\mathop{\mathit{Sh}}\nolimits}
\def\NL{\mathop{N\!L}\nolimits}
\def\CH{\mathop{\mathrm{CH}}\nolimits}
\def\proetale{{pro\text{-}\acute{e}tale}}
\def\etale{{\acute{e}tale}}
\def\QCoh{\mathit{QCoh}}
\def\Ker{\mathop{\mathrm{Ker}}}
\def\Im{\mathop{\mathrm{Im}}}
\def\Coker{\mathop{\mathrm{Coker}}}
\def\Coim{\mathop{\mathrm{Coim}}}

% Boxtimes
%
\DeclareMathSymbol{\boxtimes}{\mathbin}{AMSa}{"02}

%
% Macros for moduli stacks/spaces
%
\def\QCohstack{\mathcal{QC}\!\mathit{oh}}
\def\Cohstack{\mathcal{C}\!\mathit{oh}}
\def\Spacesstack{\mathcal{S}\!\mathit{paces}}
\def\Quotfunctor{\mathrm{Quot}}
\def\Hilbfunctor{\mathrm{Hilb}}
\def\Curvesstack{\mathcal{C}\!\mathit{urves}}
\def\Polarizedstack{\mathcal{P}\!\mathit{olarized}}
\def\Complexesstack{\mathcal{C}\!\mathit{omplexes}}
% \Pic is the operator that assigns to X its picard group, usage \Pic(X)
% \Picardstack_{X/B} denotes the Picard stack of X over B
% \Picardfunctor_{X/B} denotes the Picard functor of X over B
\def\Pic{\mathop{\mathrm{Pic}}\nolimits}
\def\Picardstack{\mathcal{P}\!\mathit{ic}}
\def\Picardfunctor{\mathrm{Pic}}
\def\Deformationcategory{\mathcal{D}\!\mathit{ef}}
\newenvironment{reference}{}{}
\newenvironment{slogan}{}{}
\newenvironment{history}{}{}
\usepackage{cleveref}
\crefname{originallemma}{original lemma}{original lemmas}
\Crefname{originallemma}{Original lemma}{Original lemmas}
\makeatletter
\newlabel{modules-lemma-stalk-module-differentials}{{08TE}{source}{}{sourcepointer-1}{}}
\newlabel{modules-lemma-stalk-module-differentials@cref}{{[originallemma][0][]08TE}{[1][0][]source}{}{}{}}
\makeatother
\crefname{originallemma}{original lemma}{original lemmas}
\Crefname{originallemma}{Original lemma}{Original lemmas}
\makeatletter
\newlabel{morphisms-lemma-base-change-differentials}{{01V0}{source}{}{sourcepointer-2}{}}
\newlabel{morphisms-lemma-base-change-differentials@cref}{{[originallemma][0][]01V0}{[1][0][]source}{}{}{}}
\makeatother
\crefname{originallemma}{original lemma}{original lemmas}
\Crefname{originallemma}{Original lemma}{Original lemmas}
\makeatletter
\newlabel{schemes-lemma-morphism-from-spec-local-ring}{{01J6}{source}{}{sourcepointer-3}{}}
\newlabel{schemes-lemma-morphism-from-spec-local-ring@cref}{{[originallemma][0][]01J6}{[1][0][]source}{}{}{}}
\makeatother
\title{The tangent space at a point is the dual of the fibre of the module of relative differentials}
\author{The Stacks Project Authors}
\date{stacks-project:varieties.tex:local-snapshot-20261011; GNU FDL 1.2 or later}
\begin{document}
\maketitle
\noindent\textit{Editorial note.} All delivered mathematical text is the author's own, copied from the pinned Stacks Project chapter file varieties.tex; no statement or proof step was rewritten, reordered or removed.\par\medskip
\noindent\textit{Editorial note.} The proof is complete as written; the only informal phrase is 'the reader sees that a to D(a) is a k-derivation', which states a verification that is immediate from the multiplicativity of phi and is reproduced as written.\par\medskip
\noindent\textit{Editorial note.} The counted spans carry 3 navigation stubs and no bibliography citation.\par\medskip
\noindent\textit{Editorial note.} The preamble is the chapter's retained minimal preamble (sources/stacks-varieties/retained-preamble.tex lines 1-77), a verbatim excerpt of the native preamble.\par\medskip
\noindent\textit{Editorial note.} The label lemma-tangent-space-cotangent-space is defined only in varieties.tex, so no cross-chapter restatement had to be weighed for this card.\par\medskip
\noindent\textit{Allowed standard prerequisites (editorial).}\begin{itemize}
\item The fibre of the relative differentials is unchanged by base change to the fibre. Morphisms, Lemma morphisms-lemma-base-change-differentials, invoked at proof lines 3038-3039. Supplies that the right hand side depends only on the scheme theoretic fibre.
\item The spectrum of a local ring represents the local behaviour of the tangent space. Schemes, Lemma schemes-lemma-morphism-from-spec-local-ring, invoked at proof lines 3040-3042. Supplies that replacing X by the spectrum of its local ring does not change T\_\{X/S,x\}.
\item The stalk of the sheaf of differentials at a point is the module of differentials of the local ring. Modules, Lemma modules-lemma-stalk-module-differentials, invoked at proof lines 3043-3044. Supplies that the same replacement does not change the fibre of Omega\_\{X/S\}.
\item Universal property of the module of differentials. Standard algebra, used at proof lines 3053-3055. Supplies the bijection between k-derivations and kappa-linear maps on the module of differentials.
\end{itemize}\medskip
\ifdefined\hypertarget\hypertarget{sourcepointer-1}{}\fi
\noindent\textit{Source reference.} Original lemma 08TE: sources/stacks-chow/tags.txt, source lines 11162--11162. Pointer outside the counted span, invoked once for context; the load-bearing argument of the counted statement is the author's proof reproduced in full in the proof span.\par\medskip
\ifdefined\hypertarget\hypertarget{sourcepointer-2}{}\fi
\noindent\textit{Source reference.} Original lemma 01V0: sources/stacks-chow/tags.txt, source lines 2429--2429. Pointer outside the counted span, invoked once for context; the load-bearing argument of the counted statement is the author's proof reproduced in full in the proof span.\par\medskip
\ifdefined\hypertarget\hypertarget{sourcepointer-3}{}\fi
\noindent\textit{Source reference.} Original lemma 01J6: sources/stacks-chow/tags.txt, source lines 2003--2003. Pointer outside the counted span, invoked once for context; the load-bearing argument of the counted statement is the author's proof reproduced in full in the proof span.\par\medskip
\begin{lemma}
\label{lemma-tangent-space-cotangent-space}
Let $f : X \to S$ be a morphism of schemes. Let $x \in X$.
There is a canonical isomorphism
$$
T_{X/S, x} = \Hom_{\mathcal{O}_{X, x}}(\Omega_{X/S, x}, \kappa(x))
$$
of vector spaces over $\kappa(x)$.
\end{lemma}

\begin{proof}
Set $\kappa = \kappa(x)$.
Given $\theta \in T_{X/S, x}$ we obtain a map
$$
\theta^*\Omega_{X/S} \to
\Omega_{\Spec(\kappa[\epsilon])/\Spec(\kappa(s))} \to
\Omega_{\Spec(\kappa[\epsilon])/\Spec(\kappa)}
$$
Taking sections we obtain an $\mathcal{O}_{X, x}$-linear map
$\xi_\theta : \Omega_{X/S, x} \to \kappa \text{d}\epsilon$, i.e.,
an element of the right hand side of the
formula of the lemma. To show that $\theta \mapsto \xi_\theta$ is
an isomorphism we can replace $S$ by $s$ and $X$ by the
scheme theoretic fibre $X_s$. Indeed, both sides of the
formula only depend on the scheme theoretic fibre;
this is clear for $T_{X/S, x}$ and for the RHS see
Morphisms, Lemma \ref{morphisms-lemma-base-change-differentials}.
We may also replace $X$ by the spectrum of $\mathcal{O}_{X, x}$
as this does not change $T_{X/S, x}$
(Schemes, Lemma \ref{schemes-lemma-morphism-from-spec-local-ring})
nor $\Omega_{X/S, x}$
(Modules, Lemma \ref{modules-lemma-stalk-module-differentials}).

\medskip\noindent
Let $(A, \mathfrak m, \kappa)$ be a local ring over a field $k$.
To finish the proof we have to show that any $A$-linear map
$\xi : \Omega_{A/k} \to \kappa$ comes from a unique $k$-algebra
map $\varphi : A \to \kappa[\epsilon]$ agreeing with the canonical
map $c : A \to \kappa$ modulo $\epsilon$. Write
$\varphi(a) = c(a) + D(a) \epsilon$
the reader sees that $a \mapsto D(a)$ is a $k$-derivation.
Using the universal property of $\Omega_{A/k}$ we see that each
$D$ corresponds to a unique $\xi$ and vice versa. This finishes the proof.
\end{proof}

\end{document}
